% Pattern: proof/attach/ref. A class that defines \ref again for its own
% typesetting, as \def\ref#1{{{\rm\@newref{#1}}}}, still writes a reference in
% the head of a proof. Seen in the gold pilot (set X).
% want theorem-1 statement/theorem 1 status=here labels=a.thm
% want theorem-2 statement/theorem 2 status=none labels=b.thm
% want proof-1 proof - proves=theorem-1
\documentclass{article}
\usepackage{amsthm}
\makeatletter
\def\@newref#1{\expandafter\@setref\csname r@#1\endcsname\@firstoftwo{#1}}
\def\ref#1{{{\rm\@newref{#1}}}}
\makeatother
\newtheorem{theorem}{Theorem}
\begin{document}
\begin{theorem}
\label{a.thm}Every compact space is bounded.
\end{theorem}
\begin{theorem}
\label{b.thm}Every finite space is compact.
\end{theorem}
We turn to the proofs.

\begin{proof}[Proof of Theorem \ref{a.thm}]
Take a finite cover.
\end{proof}
\end{document}
